\documentclass[12pt]{book}
\usepackage{amsmath}
\usepackage{amssymb}
\usepackage{geometry}
\geometry{margin=1in}

\title{Volume 09: Chemistry and Nucleus-Driven Framework \\ Book 01: Chemistry from Nuclear Geometry}
\author{James C. Harvey}
\date{December 2025}

\begin{document}
\maketitle
\tableofcontents

\chapter{Nucleus-Driven Chemistry}

\section*{Abstract}
Chemistry in SDT is driven by nuclear geometry. Bonding, reactivity, and molecular stability follow
from occlusion geometry and helical wake meshing. This book formalizes the nucleus-driven framework and
expresses chemical behavior as flux topology.

\section{Bonding as Helical Meshing}
Bond formation occurs when helical wakes align into a shared, stable flux pattern. The coupling strength
is proportional to the overlap of wake geometry and the reduction of occlusion barriers.

\section{Electronegativity as Magnetic Asymmetry}
Electronegativity is interpreted as asymmetric wake strength:
\begin{equation}
\chi \propto \left|\nabla \Pi\right|_{\text{bond scale}}.
\end{equation}
The greater the gradient at bonding distance, the stronger the attraction to shared flux.

\chapter{Reaction Pathways}

\section*{Abstract}
Reaction kinetics are pressure-barrier transitions. Activation energy is the flux cost to rearrange
occlusion geometry and re-lock a shared wake.

\section{Pressure Barrier Model}
\begin{equation}
E_a \sim \int_{\text{path}} \Delta \Pi \, d\ell.
\end{equation}

\chapter{Bond Types as Flux Topologies}

\section*{Abstract}
Bond types are topology classes of shared wake geometry. This chapter defines the primary classes in SDT
terms.

\section{Covalent Meshing}
Covalent bonding corresponds to mutual helical interlock:
\begin{equation}
\mathcal{M}_{\text{cov}} \propto \int \mathbf{B}_1 \cdot \mathbf{B}_2\, dV.
\end{equation}

\section{Ionic Alignment}
Ionic behavior is strong asymmetry in occlusion gradients:
\begin{equation}
\Delta \Pi_{\text{ion}} \gg 0 \quad \Rightarrow \quad \text{dominant directed flux}.
\end{equation}

\chapter{Polarity and Dipole Structure}

\section*{Abstract}
Polarity is the spatial offset between boundary curvature centers and wake centers. It is a geometric
dipole in the matrix.

\section{Dipole Moment}
\begin{equation}
\mathbf{p} = \sum_i q_i \mathbf{r}_i \ \rightarrow \ \sum_i \kappa_i \mathbf{r}_i,
\end{equation}
where curvature proxies replace charge.

\chapter{Catalysis and Path Steering}

\section*{Abstract}
Catalysis is path steering by altering occlusion geometry. It reshapes the pressure barrier without
changing the net flux endpoints.

\section{Barrier Reduction}
\begin{equation}
E_a' = E_a - \Delta \Pi_{\text{cat}},
\end{equation}
where $\Delta \Pi_{\text{cat}}$ is the catalyst-induced pressure reduction.

\chapter{Cross-Reference to Molecular Atlas}

\section*{Abstract}
Molecular-level details are recorded in the molecular atlas. This book provides the interpretive
framework.

\section{Atlas Boundary}
For molecule-specific geometries and numeric values, see the Molecular Structures volumes. This volume
defines the SDT chemistry grammar.
\end{document}
